\section{Introduction}

Cellular automata are discrete, grid-based mathematical models widely used to simulate complex dynamical systems through simple, localized rules \cite{cormen}. In this project, we explore the parallelization of the Forest Fire model, a canonical probabilistic cellular automaton introduced to study self-organized criticality \cite{drossel1992}. The model simulates the propagation of fire across a forested landscape, incorporating stochastic events such as spontaneous tree growth and random lightning strikes. 

The problem is computationally evaluated on a square two-dimensional grid of size $N \times N$, which establishes our total input size $n = N^2$. The state of the system evolves over $T$ discrete time steps. Because the state of any given cell in the subsequent generation depends exclusively on its current state and the states of its immediate orthogonal neighbors (the Von Neumann neighborhood), the sequential time complexity is strictly bounded by $\mathcal{O}(T \times N^2)$. For our parallel implementation, the relevant cost model considers both the distributed computational workload and the communication overhead. Specifically, the execution time is modeled as the sum of the local computation time, which scales inversely with the number of processors $p$, and the inter-process communication time required to synchronize the grid boundaries (halo exchange).

The primary objective of our work is to design, implement, and systematically analyze a topology-aware parallel version of this algorithm using the Message Passing Interface (MPI). We investigate how different domain decomposition strategies—namely 1D slices and 2D blocks—impact overall execution time, scalability, and the ratio of computation to communication.

The remainder of this report is organized as follows. Section 2 formally defines the problem and describes the sequential baseline algorithm. Section 3 presents our parallel design, highlighting the dependency analysis and the required communication patterns. In Section 4, we detail the topology-aware mappings applied to the grid. Section 5 discusses the technical specifics of our MPI implementation. Section 6 describes the hardware and methodological setup for our experiments, while Section 7 reports the empirical results regarding speedup and efficiency. Finally, Section 8 provides a critical discussion of our findings, and Section 9 concludes the report with our main takeaways.